\section{局限性与未来工作}

本章首先分析本文工作的局限性，并在此基础上概述值得进一步探索的潜在未来研究方向。

\subsection{局限性}

本文已经在基于真实描述创建 ChartRWC 数据集并实现基于混合检索的推荐方法方面取得了初步进展。然而，当前系统仍存在明显局限。系统目前仅限于10种核心图表类型，这限制了其在复杂数据探索中的适用性。被排除的图表类型，例如旭日图、平行坐标图和桑基图，对于分析层次结构、多维数据或流动数据非常重要。这一局限影响了系统在高级分析场景中的实用性。此外，系统功能仍较为基础，图表编辑器仅提供初级修改选项，也不支持小多图布局等高级可视化模式。这限制了用户进行多方面数据比较以及针对特定分析目标深度优化可视化的能力，从而降低了系统作为综合可视分析工具的潜力。

此外，当前系统的推荐和生成过程缺乏个性化。它为数据分析师和管理者提供相同图表，忽略了他们不同的需求。系统没有考虑用户角色或任务上下文，这是一个根本性局限。进一步而言，系统的视觉设计选项是预设的，无法灵活适应特定设计规范，例如视觉识别标准。这种定制化不足限制了系统在真实商业场景中的实际应用，因为这些场景通常需要量身定制的设计元素。

\subsection{未来工作}

为进一步提升本文工作的全面性和实用性，后续工作将从以下三个主要方向进行深化和拓展。

\textbf{1. 扩展图表类型覆盖范围。}
为提高图表类型表示的完整性，后续将从广度和深度两个方面系统扩展 ChartRWC 数据集。在广度方面，将引入更多来自学术界和工业界的高级复杂图表类型，以满足更广泛的分析需求。在深度方面，将开发分层图表分类体系，明确每个子类型的具体使用场景和语义，从而使模型能够生成更加准确且可解释的推荐结果。

\textbf{2. 演进为智能可视分析助手。}
为克服当前系统在支持高级可视化模式方面的局限，后续计划将系统从基础图表推荐器演进为智能可视分析助手。通过将图形语法作为底层生成引擎，系统将为小多图和分层图等视图组合方式提供坚实的理论与技术基础，尤其面向多变量分析场景。此外，后续还将增强交互式图表编辑器，使用户能够基于推荐结果进行深度设计定制。

\textbf{3. 个性化与美学考虑。}
为弥补当前个性化和美学设计方面的不足，后续将探索用户角色与偏好建模，确保推荐结果更好契合不同用户的认知模式。同时，后续将形式化已有视觉设计理论和美学原则，并将其融入生成模块，使系统能够根据预定义设计规范（如配色方案和字体样式）自动优化图表，确保输出图表不仅内容准确，而且视觉上专业并具有良好的沟通效果。

\newpage
